\documentclass[12pt]{article}
\usepackage{styles/cwilliams-standard}
% \usepackage{enumitem}
% \usepackage{parskip}

\setclass{\IPD}
\settitle{Evaluating Design In-Class Assignment}
\setauthors{Tzu-Chen Chiu, Christopher Williams}
\setheaderauthors{Chiu, Williams}

\begin{document}

\maketitlepage

\begin{aside}{Format}
    The conclusion is on page 1 and 2, and the comparison of architectures is partially on page 1 and page 2.
\end{aside}

\section{Evaluation Criteria}

\begin{itemize}
    \item \textbf{LLM RAG Usage for generated course assignments}: Using a RAG LLM model would be beneficial 
    for this feature as it would utilize the generative capability of LLM's, but restrict and finetune
    the model's database to only what is necessary - uploaded course material and specific .edu sites 
    (yes, we used an emdash of our own volition). This would also then summarily lead to us 
    using the RAG model for grading the assignments because the LLM would be pulling from the same 
    material it used to create the assignment, to grade the assignment.
    \item \textbf{Using a frontend that displays assignments}: Using a dedicated front end 
    to display the generated assignments, as well as the assignments turned in by the students to display
    the autograding of their assignments would be beneficial for ease of use and access for instructors.
    \item \textbf{Dedicated services for grading and delivery}: Having dedicated, seperate services 
    for grading assignments and ``delivering'' assignments to the students would be beneficial as 
    these services would not directly rely on another service working before hand. And, the separation 
    of these services makes maintainability much more feasible.
\end{itemize}

\section{Evaluate 6 Architectures}

\begin{table}[h!]
\centering
\caption{Architecture Model Comparison Against Design Criteria}
\renewcommand{\arraystretch}{1.6}
\begin{tabularx}{\textwidth}{|>{\bfseries}p{2.5cm}|X|X|X|c|}
\hline
\textbf{Architecture} & \textbf{LLM RAG for Assignment Generation} & \textbf{Frontend for Assignment Display} & \textbf{Dedicated Grading \& Delivery Services} & \textbf{Total Score}\\
\hline
Monolithic       & \textbf{Low} --- RAG pipeline tightly coupled; hard to scale LLM calls independently & \textbf{Medium} --- Frontend achievable but scaling UI alongside backend is inflexible & \textbf{Low} --- Grading and delivery share resources with all other logic; no clean separation & 4\\
\hline
Micro-services    & \textbf{High} --- RAG service isolates LLM logic; scales independently from other services & \textbf{High} --- Dedicated frontend service cleanly separated from backend services & \textbf{High} --- Grading and delivery naturally map to discrete, independently deployable services & 9\\
\hline
Event-Driven     & \textbf{High} --- Async events suit long-running RAG/LLM jobs well & \textbf{Medium} --- Frontend possible but event-driven adds complexity to real-time UI updates & \textbf{High} --- Grading and delivery triggered by events; fully decoupled by design & 8\\
\hline
Serverless       & \textbf{Low} --- LLM/RAG calls exceed typical function time limits; cold starts hurt UX & \textbf{Medium} --- UI can be served statically & \textbf{Medium} --- Separation is natural but execution time limits hinder grading workloads &5\\
\hline
Edge Computing   & \textbf{Low} --- RAG requires centralized model/vector DB & \textbf{Low} --- Latency optimization irrelevant for this use case & \textbf{Low} --- No benefit to pushing grading or delivery logic to edge nodes &3\\
\hline
P2P              & \textbf{Low} --- RAG needs a consistent, centralized knowledge base; P2P undermines this & \textbf{Low} --- No central server means inconsistent UI/content delivery & \textbf{Low} --- No authoritative node for grading; consistency and integrity cannot be guaranteed & 3\\
\hline
\end{tabularx}
\end{table}

\noindent\textit{Note:} For the purpose of comparison, qualitative ratings were mapped to numeric
values to compute an overall score. Each criterion was scored as \textbf{High = 3}, \textbf{Medium = 2},
and \textbf{Low = 1}, and the total score for each architecture is the sum across all three criteria.
This scoring provides a concise summary of how well each architecture aligns with the system
requirements.

\section{Best Architecture Choice}

The best architecture design choice for this system is \textbf{Microservices}, as it cleanly separates the major functions of this diverse product and enables independent scalability and maintainability, resulting in long-lasting and easy-to-understand services. This architecture aligns well with the system’s requirements for LLM-based RAG assignment generation, instructor review, and automated grading and delivery. By isolating LLM-intensive workloads from the user interface and delivery logic, microservices allow each component to scale and evolve independently without impacting overall system stability.
The primary weakness of a microservices architecture is its increased operational complexity, including service coordination, deployment, and monitoring overhead, which may be challenging for a small three-person development team. If the system were intended only as a lightweight prototype or expected to support a small user base with limited concurrency, a monolithic or serverless architecture could be a better fit due to lower setup costs. Alternatively, if future requirements emphasized large-scale asynchronous processing, such as batch grading across many courses, an event-driven architecture might be more appropriate.
To mitigate these risks, the system can adopt coarse-grained microservices initially and incrementally refine service boundaries as requirements grow. Selective use of event-based communication for long-running tasks can further improve resilience while avoiding unnecessary architectural complexity.
\end{document}